{\color{unimain}\section[(обязательное) Схема электрическая соединений рядов зажимов]{(обязательное)\\Схема электрическая соединений рядов зажимов}\label{app:a41}}
\color{black}

\newcount\figcounter
\figcounter=1
\loop
  \IfFileExists{\apppath/Приложение. Схема Э4.1/_latex/img/СЭС_\the\figcounter.pdf}{%
    \begin{figure}[H]
      \centering
      % Для первой схемы чуть меньше запаса, для остальных - максимум
      \ifnum\figcounter=1
        \def\scalefactor{0.85}%
      \else
        \def\scalefactor{0.95}%
      \fi
      
      % Убрали trim и clip, так как Python уже сделал идеальный отступ
      \includegraphics[
        width=\scalefactor\textwidth,
        height=\scalefactor\textheight,
        keepaspectratio
      ]{%
        \apppath/Приложение. Схема Э4.1/_latex/img/СЭС_\the\figcounter.pdf
      }%
      
      \caption{Схема электрическая соединений рядов зажимов (часть \the\figcounter)}%
      \label{appvo:fig_\the\figcounter}%
    \end{figure}%
    \advance\figcounter by 1
  }{%
    \figcounter=0 % файл не найден — завершаем цикл
  }%
\ifnum\figcounter>0 \repeat

\clearpage